\documentclass[a4paper,11pt]{scrartcl}

\usepackage[margin=1in]{geometry}
\usepackage[scaled]{helvet}
\usepackage[T1]{fontenc}
\usepackage[utf8]{inputenc}
\usepackage{amsmath}
\usepackage{mathptmx}
\usepackage{courier}
\usepackage{graphicx}
\usepackage{ulem}
\usepackage{bookmark}
\usepackage{paralist}
\usepackage{ngerman}
\usepackage{fancyhdr}
\usepackage{float}
\usepackage{array}
\usepackage{lipsum}


\graphicspath{ {../img/} }
\renewcommand\familydefault{\sfdefault}




\pagestyle{fancy}
\fancyhf{}
\renewcommand{\headrulewidth}{0pt}
\fancyfoot[C]{\includegraphics[width=\textwidth]{Polygon_gruen}\\ \thepage}


\rhead{\includegraphics[width=\textwidth]{LogoHeader}}
\setlength\headheight{30pt}
\setlength\footskip{15pt}

\begin{document}
\renewcommand*{\arraystretch}{1.2}
\pagenumbering{gobble}
\begin{titlepage}
    \begin{center}
        \vspace*{1cm}\Huge
        \textbf{Softwaredesign Dokumentation}\par                
        \vspace{0.5cm}\LARGE        
        Software Engineering II\par           
        \vspace{2cm}
        \includegraphics[width=0.5\textwidth]{OptimaLogo_long}\par   
        \vspace{1cm}
        \textbf{Projekttitel: BonoboBoard}\par        
        \vfill\Large   
        Jakob Hutschenreiter (1419081)\\Jiesen Wang (9839152)\\Nick Kramer (3122448)\\Patrick Küsters (2598689)\\Peter Moritz Hinkel (2783930)\par
        %\vspace{2cm}  
        %\includegraphics[width=0.5\textwidth]{Bonobo_Logo}\par        
        \vspace{2cm}
        DHBW Mannheim\\
        \today     
    \end{center}
\end{titlepage}

\section*{Änderungshistorie}
\begin{table}[h]
	\begin{tabular}{@{} p{20mm} p{25mm} p{25mm} p{75mm}}
		\textbf{Revision} & \textbf{Datum} & \textbf{Autor(en)} & \textbf{Beschreibung}\\
		1.0 & 24.01.2022 & JW & A: 1, 3 \\ 
		1.1 & 25.01.2022 & JW|PH|NK|PK & A: 2, 4, 5  \\ 
        1.2 & 26.01.2022 & JH & C: 4, A: 6 \\ 
        1.3 & 28.01.2022 & NK|PK & C: 4, 5 \\ 
	\end{tabular}
\end{table}
\noindent
Abkürzungen: Hinzugefügt/Added (A), Änderung/Changed (C), Löschung/Deleted (D)
\vspace{2cm}
\tableofcontents
\newpage
\pagenumbering{arabic}

		%------------------------------------------------------------
		%-----  -----  ----- Begin actual content -----  -----  -----
		%------------------------------------------------------------



\section{Motivation}
Um sicher zu gehen, dass alle Softwaredesignziele des BonoboBoard-Projekts richtig umgesetzt werden, werden diese in diesem Dokument festgehalten. Das Dokument erleichtert die Analyse, Planung, Implementierung und Entscheidungsfindung. Anders als im Anforderungsdokument wird nun beschrieben, wie die Ziele umgesetzt werden. 

\section{Entwicklersetup}
Verwendete Software und Libraries Versionen:
\begin{table}[H]
\begin{tabular}{|p{5cm}|p{5cm}|}
\hline
\textbf{Software/Library} & \textbf{Version} \\ \hline
	Django &  4.0.2\\ \hline
	Django Tables & 2.4.1 \\ \hline
	dj-database-url & 0.5.0 \\ \hline
	Gunicorn & 20.1.0 \\ \hline
	Docker & 20.10.12, Build e91ed57 \\ \hline
	PyCharm & 2021.3.2 \\ \hline
	VisualStudio Code & 1.64.2 \\ \hline
	Beautiful Soup & 4.8.2 \\ \hline
	Requests & 2.27.1 \\ \hline
	iCalendar & 4.0.9 \\ \hline
	Pandas & 1.4.0 \\ \hline
	lxml & 4.5.0 \\ \hline
	SQLAlchemy & 1.4.31 \\ \hline
	asgiref & 3.4.1 \\ \hline
	autopep8 & 1.6.0 \\ \hline
	Python & 3.9 \\ \hline
	pycodestyle & 2.8.0 \\ \hline
	python-decouple & 3.5 \\ \hline
	pytz & 2021.3 \\ \hline
	sqlparse & 0.4.2 \\ \hline
	toml & 0.10.2 \\ \hline
	Unipath & 1.1 \\ \hline
	whitenoise & 5.3.0 \\ \hline
\end{tabular}
\end{table}
	\subsection{Hosting bei 1Blu}
Um die Webanwendung anderen Studierenden zur Verfügung stellen zu können, haben wir uns dafür entschieden, die Anwendung auf einem virtuellen privaten Server (VPS) zu betreiben. Der VPS läuft bei dem Hosting-Anbieter 1Blu. 
Da die Anwendung mit Docker bereitgestellt wird, hält sich der Konfigurationsaufwand auf dem Server in Grenzen. Auch ein einfaches Neu-Bauen ist möglich, was im automatisierten Deployment-Prozess genutzt wird (siehe Abschnitt \ref{ci}).

\section{Verwendete Tools}
Für ein gemeinsames strukturiertes Arbeiten am Projekt werden je nach Aufgabe verschiedene Tools eingesetzt. 
	\subsection{Tools für das Projektmanagement/Organisation}
\begin{table}[H]
\begin{tabular}{|p{4cm}|p{8cm}|}
\hline
\textbf{Tool} &\textbf{Einsatzzweck} \\ \hline
Discord &  Austausch von Nachrichten, Virtuelle Meetings, Umfragen, Informationsmanagement\\ \hline
WhatsApp &  Austausch von Nachrichten\\ \hline
DropBox & Informationsmanagement \\ \hline
Google Kalender & Terminplanung   \\ \hline
Jira &  Aufgaben- und Projektmanagement, Prozessmanagement \\ \hline
Latex und Git &  Dokumente erstellen und bearbeiten \\ \hline
OneNote &  Informationsmanagement, Dokumentvorlagen erstellen und bearbeiten \\ \hline
Draw.io &  Diagramme erstellen und bearbeiten \\ \hline
Git und GitHub & Versionskontrolle und Repository-Verwaltung \\ \hline
Adobe XD & Ideen Entwicklung, Erstellung von Mockups \\ \hline
\end{tabular}
\end{table}
	\subsection{Tools für die Softwareentwicklung}
Diese sind dem separaten Dokument: \glqq Tools für die Softwareentwicklung\grqq{} zu entnehmen.

\section{Zentrale Designentscheidung}
Im Folgenden werden die zentralen Designentscheidungen aufgegriffen und erläutert.
    \subsection{Backend: SOLID-Prinzip}
Für die Interaktion mit den DHBW Webseiten bietet sich das SOLID-Prinzip an. \\
Gründe hierfür sind:
\begin{itemize}
    \item{Die Zugriffsmöglichkeiten auf die benötigten Inhalte (DHBW-Webseiten) sind Subjekt für Änderungen,
    daher müssen diese Inhalte abstrahiert werden (mittels Klassen und Klassenattributen),
    sodass das Frontend weiterhin die gleichen Schnittstellen benutzen kann.}
    \item{Die gleichen Schnittstellen für Web-Requests und das Frontend erleichtern das standardisierte Persistieren der Daten
    in Datenbanken.}
    \item{Zur Interaktion werden HTTP Anfragen benötigt. Somit bietet sich eine Oberklasse A an,
    welche diese Logik implementiert und den Unterklassen anbietet.}
    \item{Für die Webseiten Zimbra, Dualis und Moodle ist eine Authentifizierung des Endnutzers erforderlich.
    Folglich ist eine weitere Oberklasse B, welche die Oberklasse A um die Funktionalität zum Authentifizieren erweitert, vonnöten.}
    \item{Die Interaktion mit Dualis, Moodle und dem Vorlesungsplan hat als Ziel den Erhalt und die Filterung der Daten;
    Zimbra hat allerdings den Erhalt / die Filterung der Daten und das Verschicken der Daten als Funktionalität.
    Hier bietet sich es ebenfalls an, zwei Oberklassen einzuführen, welche die eben genannten Punkte unterstützen.}
\end{itemize}

    \subsection{Frontend: DRY-Prinzip}
Im Frontend erfolgt Zugriff auf die abstrakte Darstellung der \textquotedbl Datensätze\textquotedbl{} aus dem Backend und
die vorgegebenen Funktionen / Klassen von Django. Die einzige Aufgabe besteht darin, die Daten zu visualisieren und
den entsprechenden Pfaden zuzuordnen. Mit diesen Gegebenheiten wird im Frontend \textquotedbl nur\textquotedbl{}
darauf geachtet, dass keine Code-Duplizierungen auftreten, wie es das DRY-Prinzip vorgibt.

	\subsection{Automatisierter Workflow (CI-Pipeline)} \label{ci}
Damit der Test und Deployment-Prozess nicht manuell ausgeführt werden muss, nutzen wir GitHub-Actions zur Verwaltung der Continous Integration (CI).
Dabei nutzen wir zwei Workflows:
\begin{itemize}
	\item Ausführen der Unittests bei jeder Eröffnung eines Pull-Request
	\item Ausführen der Unittests und automatisches Deployment auf den VPS bei Änderungen auf dem Main-Branch
\end{itemize}
Da das automatische Deployment eine komplexere Schrittfolge beinhaltet, beschreiben wir diese genauer.
Über GitHub-Actions wird ein Ordner erzeugt, welcher alle Dateien beinhaltet, um die Produktivumgebung auf dem Server zu starten. Dieses Artefakt wird auf den Server kopiert und anstelle der alten Produktivinstanz eingesetzt. Vorher ist es notwendig, die SSL-Zertifikate zu sichern und nach dem Umstellen wieder einzusetzen.

\section{UML}
Das gesamte UML ist dem Link: \href{https://raw.githubusercontent.com/Software-Engineering-DHBW/BonoboBoard/main/documents/img/BonoboUML.drawio.svg}{BonoboBoard UML}\footnote{https://raw.githubusercontent.com/Software-Engineering-DHBW/BonoboBoard/main/documents/img/BonoboUML.drawio.svg} zu entnehmen. 

\section{Zusammenfassung}
Das BonoboBoard wird sich im Frontend am DRY-Prinzip und im Backend an dem SOLID-Prinzip orientieren.
Eine erste prototypische Version dieser Orientierung ist im UML-Diagramm dargestellt.
Das Projekt wird auf einem VPS bei 1Blu gehostet und mit Docker \textquotedbl containerized\textquotedbl.
Auf dieser \textquotedbl Containerization\textquotedbl{} baut die CI-Pipeline auf, für ein automatisches Testing
und Deployment.


		%------------------------------------------------------------
		%-----  -----  ------ End actual content ------  -----  -----
		%------------------------------------------------------------
\end{document}